% =====================================================================
%  TRABAJOS FUTUROS
% =====================================================================
\chapter*{Trabajos futuros}
\addcontentsline{toc}{chapter}{TRABAJOS FUTUROS}

\begin{notaavance}[lista sujeta a lo que arroje la campaña física]
Esta lista se cierra cuando cierre la toma de datos, el 16 de octubre de 2026.
Las cuatro líneas que siguen salen de límites ya medidos y no van a
desaparecer; lo que puede ocurrir es que las sesiones que faltan añadan alguna
más.
\end{notaavance}

Las líneas que siguen salen de límites concretos que el trabajo encontró y
documentó, no de una lista genérica de mejoras.

\textbf{Resolver la aproximación fina del vehículo.} Es la limitación que impide
cumplir la tolerancia de llegada, y está localizada. El vehículo no se mueve por
debajo de 0,40~m/s: la cadena de tracción traduce esa velocidad a cero. Hay
dos vías. Una es de hardware —calibrar los servos o cambiar la relación de
transmisión— y daría al vehículo un régimen lento del que hoy carece. La otra es
de criterio: fijar la tolerancia de llegada a un valor compatible con la
plataforma, con justificación previa y no a la vista de los resultados.

\textbf{Añadir una fuente de odometría independiente del láser.} La odometría
actual deduce el movimiento comparando barridos, y eso falla por construcción en
pasillos uniformes. El proyecto lo resolvió eligiendo un entorno con estructura,
pero esa solución limita dónde puede operar el sistema. Una unidad inercial, o
odometría visual con las cámaras que los vehículos ya llevan, daría información
de avance que no depende de la geometría del sitio.

\textbf{El ascensor como segunda vía de transición.} Hoy el cambio de piso se
hace siempre por la escalera, y el ascensor es un destino más. Admitir los dos
exige que el planificador elija entre dos puntos de transferencia por nivel,
presumiblemente el más cercano al punto anterior. El código ya nombra esa
decisión y falla explícitamente cuando encuentra dos, de modo que el punto de
extensión está identificado.

\textbf{Resolver una esquina.} Todos los recorridos medidos fueron sobre tramos
rectos. El guiado real tiene esquinas. El controlador demostró mover la dirección, pero
no se ha probado que el vehículo resuelva un cambio de rumbo grande en un pasillo
estrecho. Un intento de media vuelta en un corredor de 2,3~m no fue
posible.

\textbf{Ampliar la muestra de la campaña.} Con treinta misiones el intervalo de
confianza mide 24 puntos porcentuales. Ese ancho es de diseño y solo se reduce
con más repeticiones. Una campaña mayor permitiría además comparar las dos
condiciones experimentales, cosa que hoy no se sostiene porque sus intervalos se
solapan.

\textbf{Operar con los dos vehículos en el mismo grafo.} Las pruebas de hardware
se hicieron con un vehículo cada vez, porque los dos publican en los mismos
nombres de tópico y una orden alcanza a ambos. Separarlos por espacio de nombres
está diseñado y probado en el computador, pero no ejercitado en campo, y es
condición para la demostración del sistema completo.
